% BrandShield AI - initial draft for faculty review (detailed version)
% Compile on Overleaf with pdfLaTeX. Template: IEEEtran conference mode.

\documentclass[conference]{IEEEtran}
\IEEEoverridecommandlockouts

\usepackage{cite}
\usepackage{amsmath,amssymb,amsfonts}
\usepackage{textcomp}
\usepackage{booktabs}
\usepackage{array}
\usepackage{url}

\begin{document}

\title{BrandShield AI: A Real-Time System for Multilingual Brand Sentiment Analysis and Crisis Detection on Social Media}

\author{
\IEEEauthorblockN{Soham Borse, Tanish Bantwal, Aryan Baakle, Omsingh Rajput}
\IEEEauthorblockA{\textit{Dept. of Artificial Intelligence and Data Science} \\
\textit{K. K. Wagh Institute of Engineering Education and Research} \\
Nashik, Maharashtra, India}
\and
\IEEEauthorblockN{Dr. Darshan Medhane (Guide)}
\IEEEauthorblockA{\textit{Dept. of Artificial Intelligence and Data Science} \\
\textit{K. K. Wagh Institute of Engineering Education and Research} \\
Nashik, Maharashtra, India}
}

\maketitle

\begin{abstract}
The reputation of a brand is now largely decided by what customers say about it online. A handful of
negative comments on YouTube, Instagram or an app store can spread to thousands of people within a
few hours, and by the time a company notices, the damage is often done. Existing monitoring tools are
either expensive enterprise products or simple dashboards that report overall sentiment without
explaining its cause. They also tend to perform poorly on the Hinglish and regional-language comments
that make up a large part of Indian social media. In this paper we present BrandShield AI, a
real-time system that collects public brand mentions from multiple platforms, classifies their
sentiment using a multilingual DistilBERT model supported by a code-mixed opinion lexicon, breaks each
opinion down by aspect (such as delivery, price or service), and detects sudden negative surges using
Isolation Forest combined with a robust statistical threshold. A brand-risk score and live alerts are
presented through a web dashboard. Preliminary experiments on a hand-labelled set of 83 multilingual
comments give 78.3\% accuracy and 77.6\% macro-F1, aspect detection reaches an F1 of 0.806, and the
system processes a comment in about 9.5~ms on an ordinary laptop CPU without any GPU.
\end{abstract}

\begin{IEEEkeywords}
sentiment analysis, brand reputation, aspect-based sentiment analysis, DistilBERT, code-mixed text,
anomaly detection, Isolation Forest, real-time systems
\end{IEEEkeywords}

%=====================================================================
\section{Introduction}

\subsection{Background}
Over the last decade, social media has changed the way customers interact with brands. Earlier, a
dissatisfied customer would call a helpline or write an email that stayed private. Today the same
customer posts a comment under a YouTube video, tags the brand on Instagram, or leaves a one-star
review on the app store. These posts are public, they are read by other potential customers, and
they can be shared and amplified very quickly. As a result, online conversation has become one of
the most important signals of how a brand is actually perceived.

For businesses, this conversation is useful and risky at the same time. It is useful because it is a
free and honest source of feedback about products and services. It is risky because a small issue,
such as a batch of late deliveries or a bug in an app update, can turn into a public crisis if it is
not noticed and handled early.

\subsection{Motivation}
In our discussions with small business owners and from our own observation, we found that most of
them monitor social media manually, by scrolling through comments when they have time. This does not
scale beyond a few dozen comments a day. Commercial tools such as social listening platforms do
exist, but their pricing is aimed at large enterprises. Even these tools mainly report what
percentage of mentions are positive or negative, which is not enough to act on. A manager needs to
know \textit{why} people are unhappy and \textit{whether} the situation is getting worse.

\subsection{Challenges}
Building such a system for the Indian context involves several practical challenges:
\begin{itemize}
\item \textbf{Code-mixed language:} Comments are frequently written in Hinglish (``delivery bahut
late thi'') or in Hindi and Marathi using either Roman or Devanagari script. Most sentiment models
are trained on clean English text and do not handle this well.
\item \textbf{Mixed opinions:} A single comment often contains both praise and complaint, for
example ``food was great but delivery was late''. Assigning one label to the whole comment hides
useful information.
\item \textbf{Noisy text:} Social media text contains emojis, slang, repeated characters, hashtags,
links and a fair amount of spam.
\item \textbf{Normal fluctuations vs. real crises:} The number of negative comments rises every day
during busy hours simply because overall activity rises. A useful alert system must separate these
normal peaks from genuine problems.
\item \textbf{Limited hardware:} For the system to be affordable, it should run on a normal CPU
without requiring a GPU server.
\end{itemize}

\subsection{Contributions}
The main contributions of this work are:
\begin{enumerate}
\item A complete real-time pipeline that connects data collection, text processing, analytics and
alerting through a message queue, so that a slow platform API does not block the rest of the system.
\item A hybrid sentiment classifier that combines a multilingual DistilBERT model with a small
Hinglish, Hindi and Marathi opinion lexicon and a confidence-based rule for the neutral class.
\item Clause-level aspect-based sentiment analysis using a lexicon that the user can edit at runtime.
\item A crisis detection method that uses the \textit{share} of negative comments rather than their
raw count, combining Isolation Forest with a median/MAD based threshold.
\item A preliminary evaluation on real data, including a comparison with user-given star ratings.
\end{enumerate}

\subsection{Organization of the Paper}
Section~II reviews related work. Section~III states the problem, objectives and scope. Section~IV
describes the proposed system in detail. Section~V covers implementation, and Section~VI presents
preliminary results. Section~VII discusses limitations and future work, and Section~VIII concludes
the paper.

%=====================================================================
\section{Literature Survey}

\subsection{Transformer-Based Sentiment Classification}
The transformer architecture introduced by Vaswani et al.~\cite{b11} replaced recurrent networks with
self-attention and became the basis for most modern language models. BERT~\cite{b1} applied this
idea to pre-training on large text corpora and showed that a single pre-trained model can be
fine-tuned for many tasks, including sentiment classification, with state-of-the-art results.
However, BERT is heavy for real-time use on CPU. DistilBERT~\cite{b2} uses knowledge distillation to
reduce the model size by about 40\% and increase inference speed by about 60\%, while keeping around
97\% of BERT's language understanding performance. For a system that has to classify a continuous
stream of comments on modest hardware, this trade-off is very attractive.

Multilingual variants trained on text from many languages, such as XLM-R~\cite{b12}, allow a single
model to handle comments in different languages. Barbieri et al.~\cite{b10} introduced TweetEval, a
benchmark for social media classification tasks, and showed that models adapted to tweet-like text
perform better than general-purpose ones on short informal posts.

\subsection{Lexicon-Based and Hybrid Methods}
VADER~\cite{b3} is a rule-based sentiment analyser built specifically for social media. It uses a
human-rated word list along with rules for punctuation, capitalisation, degree modifiers and
negation. VADER is fast and needs no training, but it cannot understand context, sarcasm or words
outside its dictionary. Several recent systems combine a lexicon with a machine learning model to get
the advantages of both. In our work, the lexicon is not used as a general classifier. It is limited
to romanised Indian-language opinion words and is used only to correct the transformer where it is
known to be weak.

\subsection{Code-Mixed Sentiment Analysis}
Joshi et al.~\cite{b4} were among the first to study sentiment in Hindi-English code-mixed text. They
proposed sub-word level representations to deal with spelling variations that are common when Hindi
is written in Roman script. The SemEval-2020 Task~9~\cite{b5} on code-mixed tweets (Hinglish and
Spanglish) showed that even strong systems lose considerable accuracy on code-mixed input compared to
monolingual text. This confirms that code-mixing has to be handled explicitly in an Indian brand
monitoring tool.

\subsection{Aspect-Based Sentiment Analysis}
Aspect-based sentiment analysis (ABSA) goes one step beyond overall polarity and asks which aspect of
a product or service an opinion is about. The SemEval-2014 Task~4~\cite{b6} defined standard subtasks
for aspect term extraction and aspect polarity on restaurant and laptop reviews. Most high-performing
ABSA models are supervised and need aspect-annotated training data, which is costly to create and
specific to a domain. Since every brand has its own vocabulary, we use a lighter approach based on a
configurable aspect lexicon and clause splitting.

\subsection{Anomaly Detection and Forecasting}
Isolation Forest~\cite{b7} detects anomalies by building random trees and measuring how quickly a
point gets isolated. Anomalous points are isolated in fewer splits. The method is unsupervised, fast
and works well on low-dimensional data, which makes it suitable for detecting unusual time windows
in mention statistics. For forecasting, Prophet~\cite{b8} models a time series as a combination of
trend, seasonality and holiday effects. It is robust to missing data and its output is easy for
non-experts to interpret.

\subsection{Summary and Research Gap}
Table~\ref{tab:lit} summarises the main works reviewed. Most existing approaches address only one
part of the problem, whether that is sentiment classification, aspect extraction, or anomaly
detection. We found very few systems that combine aspect-level analysis with real-time crisis
detection, and fewer still that are designed for code-mixed Indian-language comments. BrandShield AI
aims to bring these pieces together in a single working system.

\begin{table}[t]
\caption{Summary of Related Work}
\label{tab:lit}
\centering
\footnotesize
\begin{tabular}{>{\raggedright}p{1.6cm} >{\raggedright}p{2.6cm} >{\raggedright\arraybackslash}p{3.1cm}}
\toprule
\textbf{Work} & \textbf{Approach} & \textbf{Limitation for our use} \\
\midrule
BERT~\cite{b1} & Pre-trained transformer & Heavy for CPU, English only \\
DistilBERT~\cite{b2} & Distilled BERT & Weak on code-mixed text \\
VADER~\cite{b3} & Rule-based lexicon & No context, English only \\
Joshi et al.~\cite{b4} & Sub-word LSTM & Sentiment only, no aspects \\
SemEval-2020~\cite{b5} & Code-mixed benchmark & Not a deployed system \\
SemEval-2014~\cite{b6} & Supervised ABSA & Needs labelled aspect data \\
Isolation Forest~\cite{b7} & Unsupervised anomaly detection & No notion of sentiment \\
Prophet~\cite{b8} & Time-series forecasting & Forecasts only, no alerts \\
\bottomrule
\end{tabular}
\end{table}

%=====================================================================
\section{Problem Statement, Objectives and Scope}

\subsection{Problem Statement}
To design and develop a real-time system that continuously monitors public social media conversation
about a brand, classifies sentiment in English as well as code-mixed Indian languages, identifies the
specific aspects behind customer opinions, and alerts the brand early when there is an unusual rise
in negative sentiment, while running on low-cost hardware.

\subsection{Objectives}
\begin{enumerate}
\item To collect brand-related comments automatically from YouTube, Instagram, news feeds and app
store reviews through their official APIs.
\item To clean and preprocess noisy social media text, including emojis, slang and spam.
\item To classify each comment as positive, negative or neutral using a multilingual transformer,
with additional handling for Hinglish, Hindi and Marathi.
\item To extract aspects such as delivery, price, service and food quality, and assign sentiment to
each aspect separately.
\item To detect harmful anomalies in negative sentiment and forecast short-term sentiment trends.
\item To combine these signals into a single brand-risk score and generate timely alerts.
\item To present the results on a live, easy-to-understand dashboard.
\end{enumerate}

\subsection{Scope}
The system is designed for a single brand workspace and works only with publicly available data
collected through official platform APIs. It does not scrape websites or access private messages.
The current version supports English, Hinglish, Hindi and Marathi, although the multilingual model can
classify comments in many other languages as well. Platforms such as X (Twitter) and Facebook are
outside the current scope due to API access costs.

%=====================================================================
\section{Proposed System}

\subsection{System Overview}
The system is organised into six layers that work as a streaming pipeline:
\begin{enumerate}
\item \textit{Collection layer:} fetches new comments from each platform.
\item \textit{Streaming layer:} a message queue that decouples collection from processing.
\item \textit{NLP layer:} preprocessing, sentiment classification and aspect extraction.
\item \textit{Analytics layer:} time-window aggregation, anomaly detection, forecasting and risk
scoring.
\item \textit{Alerting layer:} rule-based alerts with cool-down periods.
\item \textit{Presentation layer:} a web dashboard that updates live.
\end{enumerate}

\subsection{Data Collection}
Each platform has its own collector that follows a common interface, so new platforms can be added
without changing the rest of the system. Table~\ref{tab:sources} lists the sources currently
supported. All access is read-only and uses official APIs.

\begin{table}[t]
\caption{Data Sources}
\label{tab:sources}
\centering
\footnotesize
\begin{tabular}{l l l}
\toprule
\textbf{Platform} & \textbf{Access Method} & \textbf{Poll Interval} \\
\midrule
YouTube & YouTube Data API v3 & 300 s \\
Instagram & Instagram API (OAuth 2.0) & 180 s \\
News & Google News RSS & 180 s \\
App Store & Apple customer-review RSS & 180 s \\
\bottomrule
\end{tabular}
\end{table}

Collectors publish comments to a message bus instead of calling the classifier directly. We support
Apache Kafka~\cite{b9} for distributed deployment and an in-process asynchronous queue for running
on a single machine. Because of this decoupling, a burst of several thousand comments or a slow API
only increases the queue length and does not stop the system.

\subsection{Preprocessing}
Each comment passes through the following steps in order:
\begin{enumerate}
\item Removal of URLs, @-mentions and HTML tags.
\item Hashtag unwrapping, e.g. \#latedelivery $\rightarrow$ latedelivery.
\item Shortening of repeated characters, e.g. ``sooooo'' $\rightarrow$ ``soo''.
\item Conversion of about 30 common emojis into words (an angry face becomes ``angry''), since the
model tokenizer treats most emojis as unknown tokens.
\item Tokenisation and stop-word removal, while keeping negation words such as ``not'' and ``nahi''.
\item Lemmatisation using WordNet.
\item Language identification.
\end{enumerate}
A rule-based spam filter flags promotional and bot-like content such as follow-for-follow messages,
link-heavy posts and hashtag stuffing. Spam comments are stored but excluded from all sentiment
statistics.

\subsection{Hybrid Sentiment Classification}
The base model is a multilingual DistilBERT that outputs five ordinal classes (very negative to very
positive). We merge these into three classes by adding neighbouring probabilities. Let
$p_m = (p^+, p^0, p^-)$ be the resulting model distribution over positive, neutral and negative.

For code-mixed text we also compute a lexicon polarity $\ell \in [-1, 1]$, the average signed weight
of opinion words found from our Hinglish, Hindi and Marathi word list. A negation word within two
tokens flips the sign. If at least one lexicon word is found, the two distributions are blended:
\begin{equation}
p = (1 - \lambda)\, p_m + \lambda\, p_\ell, \qquad \lambda = 0.35
\end{equation}
where $p_\ell = \big(\max(\ell,0),\; 1 - |\ell|,\; \max(-\ell,0)\big)$.

We observed that the model often gives a weak positive or negative label to purely factual
comments. To correct this, if the top class is not neutral, no lexicon word matched, and the gap
between the two highest probabilities is small,
\begin{equation}
p_{(1)} - p_{(2)} < \delta, \qquad \delta = 0.20
\end{equation}
the comment is labelled neutral. Finally, a continuous sentiment score
\begin{equation}
s = p^+ - p^-, \qquad s \in [-1, 1]
\end{equation}
is stored with each comment and used in all later analytics. Results are cached on the normalised
text, since social media streams contain many repeated comments.

\subsection{Aspect-Based Sentiment Analysis}
Each comment is split into clauses at sentence boundaries and at contrast words such as ``but'',
``however'', ``lekin'', ``magar'' and ``par''. Each clause is then matched against an aspect lexicon
using whole-word matching. Whole-word matching avoids errors such as detecting ``late'' inside
``chocolate''. The default lexicon has eight aspects: food quality, service, delivery, price,
ambience, hygiene, app experience and packaging. Every clause that contains an aspect is classified
separately by the sentiment model, so ``food was great but delivery was late'' produces a positive
result for food quality and a negative result for delivery. Users can add or edit aspects from the
dashboard without restarting the system.

\subsection{Crisis Detection}
Comments are grouped into fixed time buckets of 15 minutes. For each bucket $t$ we compute the total
number of comments $n_t$, the number of negative comments $n^-_t$, the mean sentiment $\bar{s}_t$,
and the negative share
\begin{equation}
r_t = \frac{n^-_t}{n_t}.
\end{equation}

Two detectors run in parallel. The first is an Isolation Forest trained on recent buckets using the
feature vector $\langle n_t, n^-_t, r_t, \bar{s}_t, \Delta n^-_t \rangle$. The second is a robust
z-score computed against a trailing window of $w$ buckets using the median and the median absolute
deviation (MAD):
\begin{equation}
z_t = \frac{x_t - \mathrm{med}(x_{t-w:t-1})}{1.4826 \cdot \mathrm{MAD}(x_{t-w:t-1})}
\end{equation}
The median and MAD are used instead of the mean and standard deviation because they are not
distorted by the very spikes we are trying to detect.

A bucket is marked as a crisis only when the negative \textit{share} rises sharply and this is
confirmed by the negative \textit{count}. In our early tests, a rule based only on the raw count
flagged normal lunch and evening peaks as crises, because the count goes up whenever overall
activity goes up. Using the share as the main signal removed these false alarms.

\subsection{Forecasting and Risk Score}
Hourly series of mean sentiment, negative share and volume are forecast with Prophet~\cite{b8}. When
fewer than 36 hours of data are available, the system falls back to Holt's linear smoothing.

A brand-risk score $R \in [0, 100]$ is computed from five components, each scaled to $[0,1]$:
\begin{equation}
R = 35 f_1 + 20 f_2 + 25 f_3 + 10 f_4 + 10 f_5
\end{equation}
where $f_1$ is the recent negative share, $f_2$ the speed at which negative comments are increasing,
$f_3$ the anomaly score, $f_4$ the forecasted decline over the next six hours, and $f_5$ how much of
the negativity is concentrated in a single aspect. The score is mapped to four levels: low, medium,
high and critical. Since a high score needs several signals to agree, a single noisy signal does not
cause a false alarm.

\subsection{Alerting and Dashboard}
Alert rules cover negative spikes, risk threshold crossings, sustained negativity, aspect-specific
surges, forecasted decline, spam waves and individual negative posts going viral. Each rule has a
30-minute cool-down so that one incident produces one alert and not a flood of messages. Alerts are
pushed to the dashboard over WebSocket and can also be sent by email, Slack, Microsoft Teams or SMS.
The dashboard is a React application showing live sentiment trends, aspect breakdowns, the risk
score, forecasts and the alert history.

%=====================================================================
\section{Implementation Details}

The backend is written in Python using FastAPI, and the models run on PyTorch through the Hugging Face
Transformers library. Table~\ref{tab:stack} lists the main tools used. Development was done on a
laptop with an AMD Ryzen~7 7435HS processor and 24~GB RAM. Although the laptop has an RTX~4050 GPU,
it was intentionally not used, so that the results reflect CPU-only deployment. The complete system
uses about 1.77~GB of memory and can be hosted on a basic 2-core, 4~GB cloud virtual machine.

\begin{table}[t]
\caption{Technology Stack}
\label{tab:stack}
\centering
\footnotesize
\begin{tabular}{l l}
\toprule
\textbf{Component} & \textbf{Tool / Library} \\
\midrule
Backend API & FastAPI (Python) \\
Deep learning & PyTorch (CPU), Hugging Face Transformers \\
Sentiment model & Multilingual DistilBERT \\
Text processing & NLTK, WordNet \\
Anomaly detection & scikit-learn (Isolation Forest) \\
Forecasting & Prophet, Holt smoothing \\
Message queue & Apache Kafka / asyncio queue \\
Database & SQLite \\
Frontend & React, Chart.js \\
\bottomrule
\end{tabular}
\end{table}

%=====================================================================
\section{Preliminary Results}

\subsection{Datasets}
\textit{Annotated benchmark (83 comments):} brand comments that we labelled by hand as positive,
neutral or negative. The set intentionally includes English, Hinglish, Hindi in both scripts,
Marathi, sarcasm, mixed opinions and spam. Aspect labels are available for 47 of these comments.

\textit{Live corpus (9,984 comments):} collected by the running system from a YouTube channel (9,549
comments over 15 videos), Google News (283 articles) and the Apple App Store (150 reviews). Apart
from English, the corpus contains Indonesian, Spanish, Portuguese, Italian and Arabic comments, which
shows that multilingual support is needed in practice.

\textit{External validation (150 reviews):} App Store reviews come with a star rating given by the
reviewer. We map 1--2 stars to negative, 3 to neutral and 4--5 to positive, and use this as an
independent label.

\subsection{Sentiment Classification}
Table~\ref{tab:results} compares the transformer alone with our hybrid classifier. Accuracy stays the
same, but macro-F1 and macro-recall improve, which means the lexicon mainly helps on the harder
minority-class and code-mixed comments. We also tried a smaller distilled multilingual model
earlier, which gave only 49.4\% accuracy. Choosing the right base model turned out to be the most
important decision in the project.

\begin{table}[t]
\caption{Sentiment Classification on the Annotated Benchmark}
\label{tab:results}
\centering
\footnotesize
\begin{tabular}{l c c c}
\toprule
\textbf{Model} & \textbf{Accuracy} & \textbf{Macro-P} & \textbf{Macro-F1} \\
\midrule
DistilBERT only & 0.783 & 0.774 & 0.765 \\
Hybrid (ours) & 0.783 & 0.784 & 0.776 \\
\bottomrule
\end{tabular}
\end{table}

Table~\ref{tab:perclass} shows per-class results. No negative comment was predicted as positive, so
the system almost never reports a happy customer when the customer is actually unhappy. The neutral
class is the weakest, as it is hard to separate factual statements from mild opinions.

\begin{table}[t]
\caption{Per-Class Performance of the Hybrid Classifier}
\label{tab:perclass}
\centering
\footnotesize
\begin{tabular}{l c c c c}
\toprule
\textbf{Class} & \textbf{Precision} & \textbf{Recall} & \textbf{F1} & \textbf{Support} \\
\midrule
Negative & 0.757 & 0.800 & 0.778 & 35 \\
Neutral & 0.636 & 0.700 & 0.667 & 20 \\
Positive & 0.958 & 0.821 & 0.885 & 28 \\
\bottomrule
\end{tabular}
\end{table}

\subsection{Aspect Detection and External Validation}
On the 47 aspect-labelled comments, aspect detection achieves a precision of 0.900, recall of 0.730
and F1 of 0.806. High precision means that when an aspect is assigned, it is almost always correct;
the lower recall is due to words not yet present in the lexicon. On the 150 App Store reviews, our
predictions matched the star-based labels for 122 reviews (81.3\%). Since these labels come from the
reviewers themselves, this gives some confidence that the benchmark accuracy is not a result of our
own labelling choices.

\subsection{Speed and Early Crisis Test}
With a batch size of 32, the system classifies a comment in 9.51~ms on the laptop CPU. In an early
test where we injected a burst of negative delivery complaints into the live stream, the first alert
was raised after about 60 seconds and the risk score rose from 12 to 68.

%=====================================================================
\section{Limitations and Future Work}

Since this is an early version, there are several limitations that we plan to address:
\begin{itemize}
\item The annotated benchmark has only 83 comments labelled by our team. We plan to build a larger
dataset with multiple annotators and report inter-annotator agreement.
\item Comments with mixed opinions are usually labelled negative overall, even though the aspect-level
output is correct.
\item Some code-mixed complaints, especially about price (``bahut mehenga hai''), are still missed.
\item Aspect detection depends on the lexicon, so new vocabulary is missed until it is added. We plan
to try a few-shot or prompt-based aspect extractor.
\item The crisis test used an injected surge with a sharp start. Real crises build up more slowly, so
we plan to evaluate on documented real-world brand incidents.
\end{itemize}

%=====================================================================
\section{Conclusion}

This paper presented the design and initial implementation of BrandShield AI, a real-time system for
monitoring brand sentiment on social media. The system combines multilingual sentiment
classification, aspect-level analysis, anomaly-based crisis detection, forecasting and alerting in a
single pipeline that runs on an ordinary CPU. Early results are encouraging, with 78.3\% accuracy on a
multilingual benchmark, 0.806 F1 for aspect detection and 81.3\% agreement with independent star
ratings. The next stage of the work will focus on a larger evaluation dataset, better handling of
code-mixed text and testing on real brand crises.

\section*{Acknowledgment}
We thank our guide, Dr. Darshan Medhane, for his guidance and support, and the Department of
Artificial Intelligence and Data Science, K. K. Wagh Institute of Engineering Education and Research,
Nashik, for providing the necessary facilities.

\begin{thebibliography}{00}
\bibitem{b1} J. Devlin, M.-W. Chang, K. Lee and K. Toutanova, ``BERT: Pre-training of deep
bidirectional transformers for language understanding,'' in \textit{Proc. NAACL-HLT}, 2019,
pp. 4171--4186.
\bibitem{b2} V. Sanh, L. Debut, J. Chaumond and T. Wolf, ``DistilBERT, a distilled version of BERT:
smaller, faster, cheaper and lighter,'' \textit{arXiv:1910.01108}, 2019.
\bibitem{b3} C. J. Hutto and E. Gilbert, ``VADER: A parsimonious rule-based model for sentiment
analysis of social media text,'' in \textit{Proc. 8th Int. AAAI Conf. Weblogs and Social Media
(ICWSM)}, 2014.
\bibitem{b4} A. Joshi, A. Prabhu, M. Shrivastava and V. Varma, ``Towards sub-word level compositions
for sentiment analysis of Hindi-English code mixed text,'' in \textit{Proc. COLING}, 2016,
pp. 2482--2491.
\bibitem{b5} P. Patwa et al., ``SemEval-2020 Task 9: Overview of sentiment analysis of code-mixed
tweets,'' in \textit{Proc. 14th Workshop on Semantic Evaluation (SemEval)}, 2020, pp. 774--790.
\bibitem{b6} M. Pontiki, D. Galanis, J. Pavlopoulos, H. Papageorgiou, I. Androutsopoulos and
S. Manandhar, ``SemEval-2014 Task 4: Aspect based sentiment analysis,'' in \textit{Proc. 8th Int.
Workshop on Semantic Evaluation (SemEval)}, 2014, pp. 27--35.
\bibitem{b7} F. T. Liu, K. M. Ting and Z.-H. Zhou, ``Isolation Forest,'' in \textit{Proc. 8th IEEE
Int. Conf. Data Mining (ICDM)}, 2008, pp. 413--422.
\bibitem{b8} S. J. Taylor and B. Letham, ``Forecasting at scale,'' \textit{The American
Statistician}, vol. 72, no. 1, pp. 37--45, 2018.
\bibitem{b9} J. Kreps, N. Narkhede and J. Rao, ``Kafka: A distributed messaging system for log
processing,'' in \textit{Proc. NetDB Workshop}, 2011.
\bibitem{b10} F. Barbieri, J. Camacho-Collados, L. Espinosa-Anke and L. Neves, ``TweetEval: Unified
benchmark and comparative evaluation for tweet classification,'' in \textit{Findings of EMNLP},
2020, pp. 1644--1650.
\bibitem{b11} A. Vaswani et al., ``Attention is all you need,'' in \textit{Proc. Advances in Neural
Information Processing Systems (NeurIPS)}, 2017, pp. 5998--6008.
\bibitem{b12} A. Conneau et al., ``Unsupervised cross-lingual representation learning at scale,'' in
\textit{Proc. ACL}, 2020, pp. 8440--8451.
\end{thebibliography}

\end{document}
